\documentclass[10.5pt,a4paper]{scrartcl}
\usepackage[english]{babel}
\usepackage{fontspec}
\IfFontExistsTF{Libertinus Serif}{\setmainfont{Libertinus Serif}\setsansfont{Libertinus Sans}[Scale=0.94]}{\setmainfont{Latin Modern Roman}\setsansfont{Latin Modern Sans}}
\usepackage{microtype}
\usepackage[a4paper,margin=20mm]{geometry}
\usepackage{multicol}
\usepackage{xcolor}
\usepackage{lettrine}
\usepackage{graphicx}
\usepackage[hidelinks]{hyperref}
\definecolor{Accent}{HTML}{A14D32}
\definecolor{Muted}{HTML}{666666}
\setkomafont{disposition}{\sffamily\bfseries}
\setlength{\parindent}{0pt}
\setlength{\parskip}{0.35\baselineskip}
\begin{document}
{\small\sffamily\bfseries\color{Accent} EDITORIAL / ISSUE 01\par}
{\fontsize{32}{36}\selectfont\bfseries Magazine Editorial\par}
{\large\color{Muted} A disciplined multi-column reference system.\par}
\vspace{1.2em}\hrule\vspace{1em}
\begin{multicols}{2}
\lettrine[lines=3]{E}{ditorial} layouts benefit from a controlled grid, strong opening hierarchy,
and deliberate variation rather than arbitrary decoration. Captions, pull quotes, and images should
support the narrative rhythm.
\section*{Grid and rhythm}
Use columns when scanning and visual pacing improve. Maintain a stable baseline and consistent gaps.
\columnbreak
{\Large\sffamily\bfseries\color{Accent} ``A pull quote earns its scale by clarifying the article's central idea.''\par}
\vspace{1em}
\section*{Image logic}
Large images establish hierarchy. Smaller images provide evidence or detail. Avoid galleries without narrative purpose.
\end{multicols}
\end{document}
